\section{Discussion}\label{sec:discussion}

\subsection{Summary of findings}
\Cref{tab:findings} lists every result with its provenance and evidence. One reproduces the paper; the
others are ours, either audits of the paper's claims, measurements, or new analysis. Where a known
technique is applied to this scheme (nonce reuse, hedged nonces, bad-signature identification), only the
application is claimed.

\begin{table}[ht]
\centering\small
\caption{All results, with provenance and evidence.}\label{tab:findings}
\begin{tabularx}{\linewidth}{>{\raggedright\arraybackslash}X>{\raggedright\arraybackslash}p{2.7cm}>{\raggedright\arraybackslash}p{4.1cm}}
\toprule
Result & Kind & Evidence \\
\midrule
The paper's malicious-KGC attack forges Verma signatures & Reproduced & 8/8 \\
The repair stops it, and the $T_i$ binding is the change that does & Reproduced + ablation & 0/8; ablation 8/8 \\
The repaired scheme costs about twice its printed row & Audit & $\times1.99$ at $n=100$; timing 2.12--2.37 \\
Anyone with a public key can forge Verma's scheme; Table II is wrong & Ours & 8/8, 8/8 in aggregate \\
Hashing $T$ and hashing $R$ each block their own attack & Ours & $2\times2$, 8 trials per cell \\
Best-effort optimisation of both schemes does not close the gap & Experiment & $\times2.37\to\times2.40$ \\
Pre-verifying a batch is cheap, so one bad signature need not sink it & Ours (applied) & $1.7\times$ per signature \\
A repeated nonce reveals every secret: 3 signatures, or 2 for the KGC & Ours (applied) & 8/8 \\
The replay and forward-security claims do not hold & Ours & both tested \\
\bottomrule
\end{tabularx}
\end{table}

\subsection{What the paper gets right}
Our findings should not be read as a rejection of the paper. Its diagnosis of Verma's scheme is correct
and reproduces exactly. Its repair resisted every forgery we tried, the paper's own attack and our key-only
attack alike, and both of its hash bindings turn out to be necessary. The aggregator uses no secret
material, which a test confirms; that aggregate security then follows is the paper's argument, which we
did not audit. Nothing we found contradicts the paper's unforgeability theorem. Our findings concern the
efficiency comparison, deployment, and the additional properties claimed in Section~V-C.

\subsection{Threats to validity}
\begin{itemize}
\item \textbf{The proof.} We did not audit the paper's EUF-CMA proof (Theorem~1 and its lemmas). Failed
      attacks are evidence about those attacks only, not about security.
\item \textbf{Verma's scheme.} We implemented it from the paper's description (Section~IV-A), not from
      Verma's original article. Our key-only forgery and the contradiction with Table~II therefore concern
      the scheme as the paper specifies and rates it.
\item \textbf{Timing.} Absolute times include Python and \code{ctypes} overhead, and our curves are not
      the paper's. We therefore report ratios between schemes on the same machine, curve and library,
      supported by exact operation counts that are identical on both libraries. Timings come from a laptop
      and one cloud machine; we did not measure on sensor hardware.
\item \textbf{Simulated faults.} The stuck generator of \cref{sec:nonce} and the faulty signatures of
      \cref{sec:agg} are simulated in software, with independent faults at fixed rates.
\item \textbf{Our own errors.} Twice our optimisation comparison was unequal (\cref{sec:opt}); both errors
      were caught and corrected, and tests now pin the corrected behaviour.
\item \textbf{Novelty.} The mechanism of our key-only forgery is textbook Schnorr hygiene; we did not
      survey the literature exhaustively and claim only that the paper does not discuss it.
\end{itemize}
